% NAF 5166, batch 1 — Gallica views 1–20.
% Pass: Opus 5.5 (claude-opus-5-5), 2026-09-23 — first pass, unchecked against the views by a human.
%
% What the views are. Greyscale scans, portrait, about 4700 × 6430, one page
% per view. The leaves of the dossier are smaller than the modern guard
% leaves and sit on them, hinged on stubs, with the ragged edges of old
% paper; the edge of the facing guard shows at the left. Every view was
% first looked at as a 1000-px thumbnail (views 3–12 and 14 also at
% stretched contrast, to rule out faint writing); the written views were
% read as four full-width bands at 2400 px (12 % overlap), with tight crops
% at 1200–2200 px for struck words, the dates, the numbers and every
% formula. Read only from the local mirror. \page{N} is always the view.
%
% What the twenty views hold.
%
%   v. 1      Outer board, marbled paper. Not transcribed.
%   v. 2      Inside of the upper board, with the printed oval shelf label
%             « FR. / NOUV. ACQ. / 5166 ». Not transcribed.
%   v. 3–11   Blank modern guard leaves, rectos and versos (checked at
%             contrast). Not transcribed.
%   v. 12     Blank; only the reversed show-through of the title on v. 13.
%             Not transcribed.
%   v. 13     Modern title leaf, gothic calligraphy, « Papiers / de /
%             Lagrange », with the library's certificate « Volume de 26
%             Feuillets / plus les Feuillets A préliminaire & 6 bis. / 27
%             Septembre 1888. » and « Libri 1856 ». Given as a \note{} only,
%             for what it says of the volume.
%   v. 14     Verso of the title leaf, blank. Not transcribed.
%   v. 15     Leaf « A. » (the certificate's « Feuillet A préliminaire »), an
%             old wrapper: « Lagrange. » in large script; a quick
%             nineteenth-century note « Lagrange / mémoires et notes
%             autographes / avec des notes etc de / Mlle Sophie Germain »;
%             « N. a. fr. 5166. », « Libri. No. 1856. », a pencilled « 1856 »,
%             « R. c. 8070 (91) ». Given as a \note{} only: none of it is
%             Lagrange's, and the claim that the volume holds Germain's notes
%             is the wrapper's, not checked here.
%   v. 16     Verso of leaf A, blank (show-through, ghost of a stamp). Not
%             transcribed.
%   v. 17     PIECE 1 begins (leaf « 1 »). A memoir in French, headed
%             « Eclaircissement d'une difficulté singuliere qui se rencontre
%             dans le calcul de l'attraction des sphéroïdes très peu
%             differens de la sphere / Par J. L. Lagrange ». Introduction:
%             d'Alembert first computed the attraction of non-elliptic
%             spheroids little different from the sphere (vols 2 and 3 of his
%             Recherches sur le système du monde); Laplace treated it more
%             generally (Mémoires de l'Académie des Sciences for 1782, and
%             vol. 2 of the Mécanique céleste), on a « beau théoreme » that
%             gives rise to a « difficulté singuliere » nobody has remarked.
%             Article 1 starts setting the notation (r, θ, ϖ, after the
%             Mémoires de 1782 p. 134 and the Mécanique céleste) and breaks
%             off at the foot of the recto in mid-sentence.
%   v. 18     Verso of leaf 1, blank (show-through, bled ink). Not
%             transcribed.
%   v. 19     PIECE 1 continued (leaf « 2 »): R, θ', ϖ' for a molecule, the
%             distance √(r² − 2rRμ + R²) with μ = cos θ cos θ' + sin θ sin θ'
%             cos(ϖ − ϖ'), the mass R² dR dθ' dϖ' sin θ', the potential V as a
%             triple integral with its limits, and the three attractions
%             dV/dr, dV/r dθ, dV/(r dϖ sin θ). Stops at « celui-ci. » about
%             two-thirds down the page; the rest of the page is blank.
%   v. 20     Verso of leaf 2, blank (show-through). Not transcribed.
%
%   Where the piece ends. The text on v. 19 closes a sentence and is
%   followed by blank paper, and v. 20 is blank: whether the memoir goes on
%   at v. 21 (leaf 3?) cannot be told from this batch. The circled note on
%   v. 17, « 13 feuillets / le 6me double », counts a dossier of thirteen
%   leaves beginning here, which suggests, without showing, that it does.
%
%   Dating. \dating{} is the catalogue's. The memoir cites vol. 2 of the
%   Mécanique céleste; that it is therefore later than that volume is this
%   pass's inference, and is not written into the file.
%
%   The hand. One quick, sloping French cursive of about 1800, fairly
%   regular, with heavy ink and several deletions made in the course of
%   writing (v. 17, 19): a draft rather than a fair copy. The pass does not
%   say whose hand it is; the title names Lagrange as the author and the
%   wrapper (v. 15) calls the papers « autographes ». Accents are sparse and
%   often reduced to a dot; they are given where they can be seen. No long s.
%   Spelling as written (« avoit », « differens », « paroît »?).
%
%   Notation, fixed for the batch. The writer's cursive d, looped like a ∂,
%   is the same letter in words (« de », « dans ») and in differentials, and
%   is set d. The angle written as a π with a curled top bar is set ϖ
%   (\varpi). The radical is a large V followed by a parenthesis, set
%   \sqrt{( )} with the parenthesis kept. The triple integral is ∫ with a
%   superscript 3, set \int^{3}. Squares are written with a small raised 2.
%
%   Numbers on the leaves, and which kind.
%   - « A. » (v. 15) and « 1 » (v. 17): ink, top right, in a posed hand like
%     a printer's figure. They match the certificate's series (26 leaves
%     plus A and 6 bis), so they behave as the library's foliation, but are
%     in ink; following the NAF 5161 practice no \folio{} is written.
%   - « 2 » (v. 19): ink, top right, but a cursive figure in an ink like the
%     text's, not the posed figure of « 1 ». It may be the same series or the
%     writer's own numbering of his sheets; the pass cannot tell, and writes
%     no \folio{}.
%   - « 13 feuillets / le 6me double » (v. 17), circled: an archivist's or
%     owner's count of the dossier.
%   - « 1856 » (v. 13 « Libri 1856 », v. 15 « Libri. No. 1856. » and a
%     pencilled « 1856 »): given with « No. » on v. 15, so a Libri number,
%     not a year — an inference of this pass.
%   - « R. c. 8070 (91) » (v. 15, 17), « ACQ. 8070 (LIBRI) » (stamp, v. 17):
%     the library's acquisition marks.
%   - In the text: « 1782 » and « p. 134 » (the Mémoires cited), « 360° »,
%     « 180° », and the article number « 1. » (v. 17), which may be « I ».
%
%   Apparatus. Seven \ill{}, all inside \struck{} (deleted words or lines,
%   v. 17 ×2, v. 19 ×5); six \uncertain{}.

% Coordinator's reconciliation (2026-09-23), after all four batches:
%   the ink figures « A. », « 1 » (v. 15, 17) and « 2 » (v. 19) are one
%   series with « 3 » … « 11 » and « 6 bis » (v. 21–39), 12 … 22 (v. 41–60)
%   and 23 … 26 (v. 62–68): the library's count of leaves, as the
%   certificate of v. 13 announces it (26 leaves, plus A and 6 bis). The
%   question left open above for « 2 » is settled so; still no \folio{},
%   the series being in ink. The memoir begun here runs on through v. 39
%   and, by its subject, notation and hand, to v. 41 (leaf 12).

\documentclass[11pt,a4paper]{article}
% The preamble shared by every transcription of the mathematicians' manuscripts in Gallica.
%
% It rests on one idea: **uncertainty compounds, it does not resolve.** A
% transcription that smooths over a doubtful word has destroyed the only thing
% it was made for — a reader must be able to tell, at every point, what was on
% the page and what was supplied. The apparatus macros below are therefore the
% critical apparatus, not decoration.
%
% The same commands are recognised by `scripts/render.mjs`, which produces the
% site's reading view, and by `scripts/tei.mjs`. A macro added here must be
% added there too, or rendering fails loudly — which is the wanted behaviour.
%
% Comments here are English, like the rest of the repository. The documents
% this preamble sets are French, and that is deliberate: see the skills.

\ProvidesPackage{germain}[2026/09/15 Transcriptions of mathematicians' manuscripts in Gallica]

\usepackage{amsmath,amssymb,amsthm}
\usepackage{mathrsfs}
% Kept from the parent preamble so the renderer's diagram support stays
% available; these manuscripts draw no commutative diagrams, and nothing here uses it.
\usepackage{tikz-cd}
\usepackage[english,french]{babel}
\usepackage{fontspec}

% --- Greek ------------------------------------------------------------------
% The text face stays fontspec's default, Latin Modern, which has no Greek:
% XeTeX drops every glyph it lacks with a line in the log and nothing on the
% page, and the Greek words the manuscripts quote vanished from the PDFs. So
% the Greek and Greek Extended blocks get a XeTeX character class of their own,
% and entering or leaving a run of it switches the family to CMU Serif — the
% same Computer Modern design, with polytonic Greek — and back. Only the
% family changes: a Greek word in \textit or \textbf keeps its shape and series.
% The transitions to and from every other class carry the tokens that class
% has towards an ordinary letter, so French spacing before « ; » or inside
% « » still holds after a Greek word. No grouping, so no brace or \endgroup
% can come out unbalanced; maths is left alone, since XeTeX inserts nothing
% there. Kept identical in `exercices.sty`.
\makeatletter
\newfontfamily\germainGreekfont{cmun}[
  Extension=.otf, NFSSFamily=germaingreek,
  UprightFont=*rm, ItalicFont=*ti, SlantedFont=*sl,
  BoldFont=*bx, BoldItalicFont=*bi, BoldSlantedFont=*bl]
\newXeTeXintercharclass\germain@greek
\def\germain@greekrange#1#2{%
  \@tempcnta=#1\relax
  \loop
    \XeTeXcharclass\@tempcnta=\germain@greek
    \ifnum\@tempcnta<#2\advance\@tempcnta\@ne\repeat}
\germain@greekrange{"0370}{"03FF}
\germain@greekrange{"1F00}{"1FFF}
\def\germain@greekprev{\rmdefault}
\def\germain@greekin{%
  \edef\germain@greekprev{\f@family}\fontfamily{germaingreek}\selectfont}
\def\germain@greekout{\fontfamily{\germain@greekprev}\selectfont}
\def\germain@greekjoin{%
  \XeTeXinterchartokenstate=\@ne
  \@tempcnta=\z@
  \loop
    \ifnum\@tempcnta=\germain@greek\else
      \XeTeXinterchartoks\germain@greek\@tempcnta=\expandafter{%
        \expandafter\germain@greekout\the\XeTeXinterchartoks\z@\@tempcnta}%
      \XeTeXinterchartoks\@tempcnta\germain@greek=\expandafter{%
        \the\XeTeXinterchartoks\@tempcnta\z@\germain@greekin}%
    \fi
    \ifnum\@tempcnta<4095\advance\@tempcnta\@ne\repeat}
% babel-french sets its own classes, and saves and restores the interchar
% state, each time a language is selected: join again after it does.
\addto\extrasfrench{\germain@greekjoin}
\addto\extrasenglish{\germain@greekjoin}
\AtBeginDocument{\germain@greekjoin}
\makeatother

\usepackage{geometry}
\geometry{a4paper,margin=25mm,marginparwidth=28mm}
\usepackage{marginnote}
\usepackage{xcolor}
\usepackage[normalem]{ulem}
\setlength{\emergencystretch}{3em}
\usepackage{hyperref}
\hypersetup{colorlinks=true,linkcolor=black,urlcolor=black,pdfborder={0 0 0}}

% --- Batch metadata ---------------------------------------------------------
% Taken as they stand by the rendering script, which shows them at the head of
% the reading view. They fix what the file claims to cover. The macro names are
% the parent project's, so that the scripts stay shared; \folder here means the
% volume's slug — `fr-9115` — and \pages counts Gallica views.

\newcommand{\folder}[1]{\gdef\grothFolder{#1}}
\newcommand{\batch}[1]{\gdef\grothBatch{#1}}
\newcommand{\pages}[2]{\gdef\grothFirst{#1}\gdef\grothLast{#2}}
\newcommand{\foldertitle}[1]{\gdef\grothFoldertitle{#1}}

% The holder's own shelfmark, printed as they write it: « Français 9115 ».
\newcommand{\shelfmark}[1]{\gdef\grothShelfmark{#1}}

% The dating, copied verbatim from the holder's catalogue — never inferred
% here. « XIXe siècle » is the BnF's; a pass that dates a leaf from its content
% says so in a \note{}, as its own inference.
\newcommand{\dating}[1]{\gdef\grothDating{#1}}

% The status notice, printed in the title block and repeated as a diagonal
% watermark on every page of the PDF. The manuscripts are in the public
% domain; what this line declares is not a right but a quality — an
% unverified first pass by a machine. The skills write the line; a file
% without it gets no watermark, so leaving it out is visible in the output.
\newcommand{\watermark}[1]{\gdef\grothWatermark{#1}}

\gdef\grothFolder{?}\gdef\grothBatch{}\gdef\grothFirst{?}\gdef\grothLast{?}%
\gdef\grothFoldertitle{}\gdef\grothShelfmark{}\gdef\grothDating{}\gdef\grothWatermark{}

% --- The résumé -------------------------------------------------------------
\newenvironment{resume}
  {\small\setlength{\parskip}{0.45em}}
  {\par\medskip\hrule\medskip}

% --- Keywords ---------------------------------------------------------------
% In English, closing the modernised reading's résumé: the modern vocabulary
% under which to search for what the leaves construct. The manifest extracts
% them to tag the volume — this line is the single source of the tags.
\newcommand{\keywords}[1]{\par\smallskip\noindent
  {\footnotesize\textsc{Keywords} --- \textit{#1}}\par}

% --- The critical apparatus -------------------------------------------------

% The Gallica view number, in the margin. This is the anchor that lets a
% disagreement over a formula be settled by looking at *one* image rather than
% a volume of 750. The site uses it to turn the facsimile as the transcription
% is scrolled. A view is one scanned image: usually one side of a leaf.
\newcommand{\page}[1]{%
  \marginnote{\footnotesize\color{gray!70}\textsf{v.\,#1}}%
  \label{page:#1}\ignorespaces}

% The BnF's pencilled foliation, where the leaf carries it and a pass can read
% it: « 198r ». This is what the literature cites — « ff. 198r–208v » — and
% the one line that turns a citation into a view. Recorded, never inferred:
% a folio a pass cannot read is not written.
\newcommand{\folio}[1]{%
  \marginnote{\footnotesize\color{gray!50}\textsf{f.\,#1}}\ignorespaces}

% The modernised reading's coarser counterpart to \page{}: which views a
% section is drawn from.
\newcommand{\pagerange}[2]{%
  \marginnote{\footnotesize\color{gray!70}\textsf{v.\,#1--#2}}%
  \ignorespaces}

% Illegible: never guessed.
\newcommand{\ill}{\textcolor{red!60!black}{[\,\ldots\,]}}

% A doubtful reading: the word is given, and flagged as uncertain.
\newcommand{\uncertain}[1]{\uline{#1}}

% An editorial addition — a missing letter, an implied word.
\newcommand{\add}[1]{\textcolor{blue!50!black}{[#1]}}

% Struck out by the writer. What was crossed out is part of the document.
\newcommand{\struck}[1]{\sout{#1}}

% The transcriber's note, on the state of the page or the mathematics.
\newcommand{\note}[1]{\par\smallskip{\small\color{gray!60!black}\textit{[#1]}}\par\smallskip}

% A marginal note of hers, distinct from the body of the text.
\newcommand{\marginal}[1]{\par\smallskip\noindent\hspace{1em}%
  {\small\color{gray!60!black}#1}\par\smallskip}

% --- Title ------------------------------------------------------------------
% The title block is French, like the documents themselves.

\AddToHook{shipout/background}{%
  \ifx\grothWatermark\empty\else
    \begin{tikzpicture}[remember picture, overlay]
      \node[rotate=52, scale=3.4, align=center, text=gray!18,
            font=\sffamily\bfseries]
        at (current page.center) {\grothWatermark};
    \end{tikzpicture}%
  \fi}

\AtBeginDocument{%
  \begin{center}
    {\large\bfseries Manuscrits de mathématiciens (Gallica) ---
      \ifx\grothShelfmark\empty\grothFolder\else\grothShelfmark\fi}\\[2pt]
    {\itshape\grothFoldertitle}\\[4pt]
    {\small vues \grothFirst--\grothLast
      \ifx\grothBatch\empty\else\ (lot \grothBatch)\fi}\\[2pt]
    \ifx\grothDating\empty\else
      {\small\color{gray!60!black}Datation du catalogue : \grothDating}\\[2pt]
    \fi
    \ifx\grothWatermark\empty\else
      {\small\bfseries\color{red!45!gray}\grothWatermark}\\[2pt]
    \fi
    {\footnotesize\color{gray!60!black}Transcription établie d'après les images IIIF de Gallica.
     Source gallica.bnf.fr / Bibliothèque nationale de France.\\
     Les lectures incertaines sont soulignées, les illisibles notées [\,\ldots\,],
     les ajouts entre crochets ; \textsf{v.} numérote les vues de Gallica,
     \textsf{f.} la foliotation au crayon de la BnF.}
  \end{center}
  \vspace{1em}}

\endinput


\folder{naf-5166}
\shelfmark{NAF 5166}
\batch{1}
\pages{1}{20}
\dating{1701-1900}
\watermark{Lecture automatique\\non vérifiée}
\foldertitle{Papiers du géomètre J.-L. LAGRANGE (1813).}

\begin{document}

\page{13}
\note{Feuillet de titre moderne, d'une main calligraphiée du XIX\textsuperscript{e} siècle, en lettres gothiques : « Papiers / de / Lagrange ». En haut à droite, dans une accolade, « Fr. nouv. acq. 5166 ». Au-dessous du titre, le certificat de la bibliothèque : « Volume de 26 Feuillets / plus les Feuillets A préliminaire \& 6 \textsuperscript{bis}. / 27 Septembre 1888. » ; en bas à gauche, de la même main, « Libri 1856 ». Ce feuillet n'est pas de Lagrange et n'est donné ici que pour ce qu'il dit du volume : vingt-six feuillets numérotés, un feuillet préliminaire coté « A » et un feuillet « 6 bis ».}

\page{15}
\note{Feuillet préliminaire, de papier ancien, monté sur onglet. En haut à droite, à l'encre, « A. » : c'est le « Feuillet A préliminaire » du certificat de la vue 13 ; aucun folio n'est porté, voir l'en-tête. À gauche, en grande cursive soulignée, « Lagrange. » ; à côté, au crayon et de biais, « 1856 ». Au milieu de la page, d'une cursive rapide du XIX\textsuperscript{e} siècle : « Lagrange / \uncertain{mémoires} et notes \emph{autographes} / avec des notes \uncertain{etc} de / M\textsuperscript{lle} Sophie Germain » (« autographes » souligné). Plus bas, « N. a. fr. 5166. » ; puis « Libri. N\textsuperscript{o}. 1856. » ; au pied, « R. c. 8070 (91) » souligné d'une accolade. Aucune de ces mentions n'est de Lagrange : ce sont des notes de possesseur et de bibliothèque, données pour ce qu'elles disent du volume, et notamment qu'il contiendrait des notes de Sophie Germain. « mémoires » : la fin du mot est un simple trait, le pluriel n'est pas sûr ; « etc » : trois lettres rapides, lues avec doute.}

\page{17}
\note{Feuillet de papier ancien, le bord supérieur déchiré, monté sur onglet. En haut à droite, à l'encre, un « 1 » de la même main posée que le « A. » de la vue 15 ; aucun folio n'est porté, voir l'en-tête. À gauche de ce chiffre, la cote « FR. nouv. acq. 5166 » soulignée d'une accolade. Dans la marge droite, dans un cercle tracé à l'encre, d'une autre main que le texte : « 13 feuillets / le 6\textsuperscript{me} double » — un compte du dossier, qui répond au feuillet « 6 bis » du certificat (vue 13). Au milieu de la page, le cachet rond « BIBLIOTHÈQUE NATIONALE / R.F. / MANUSCRITS », posé sur les mots « très remarquable » ; au pied, « R. c. 8070 (91) » souligné d'une accolade et le cachet ovale « ACQ. 8070 (LIBRI) ». Le texte n'occupe que le recto ; le verso (vue 18) est blanc.}

\section*{Eclaircissement d'une difficulté singuliere}

Eclaircissement d'une difficulté singuliere qui se rencontre dans le calcul de l'attraction des sphéroïdes très peu differens de la sphere

Par J. L. Lagrange
\note{Titre en trois lignes, puis la ligne d'auteur, en retrait ; même main et même encre que le texte. Les accents sont notés tels qu'ils se lisent ; plusieurs manquent ou se confondent avec un point.}

D'alembert est le premier qui ait donné le calcul de l'attraction des sphéroïdes non elliptiques mais peu differens de la sphere, dans le second, et \struck{\ill{}} troisieme volumes de ses Recherches sur le systeme du monde. M. Laplace a traité ensuite cette matiere d'une maniere nouvelle et plus generale qu'on ne l'avoit fait, dans les Memoires de l'Academie des Sciences de 1782, et dans le second volume de sa Mecanique celeste. Sa théorie est fondée sur un beau théoreme très remarquable par sa simplicité autant que par sa généralité ; mais ce theoreme donne lieu à une difficulté singuliere qui \uncertain{paroît} n'avoir encore été remarquée par personne et qui merite d'être examinée.
\note{« et … troisieme » : entre « et » et « troisieme », un mot biffé d'un trait épais, et au-dessus un second mot, biffé lui aussi ; ni l'un ni l'autre ne se lit. « 1782 » : le deuxième chiffre est un 7 barré, lu ainsi ; même millésime à l'article 1. « paroît » : la voyelle après « par » peut être un \emph{o} ou un \emph{a}.}

1. Soit comme dans les Memoires de 1782 p. 134 et dans la Mecanique celeste (tome 2) $r$ la distance du point attiré au point pris pour centre du sphéroïde, $\theta$ l'angle que le rayon $r$ fait avec un axe fixe passant par le centre, $\varpi$ l'angle que \struck{\ill{}} le plan passant par cet axe et par le rayon $r$ fait
\note{« 1. » : le numéro d'article est un trait vertical à empattements, qui peut être « 1 » ou « I ». « (tome 2) » est écrit dans la marge droite et appelé par un signe d'insertion après « Mecanique celeste ». « $\varpi$ » : un $\pi$ au trait supérieur recourbé, lu comme le $\varpi$ de l'époque. Après « l'angle que », un mot court biffé, commençant par une haste, illisible. La phrase continue en tête de la vue 19.}

\page{19}
\note{Deuxième feuillet du mémoire, de même papier, le bord supérieur déchiré. En haut à droite, à l'encre, « 2 », d'un trait cursif et non de la main posée du « 1 » de la vue 17 ; voir l'en-tête. Pas de cachet ni de cote. La page commence au milieu de la phrase laissée à la fin de la vue 17.}

fait avec un plan invariable passant par le même axe ; soit de plus $R$ la distance d'une molecule du sphéroïde au centre, et soient $\theta'$, $\varpi'$ ce que deviennent les angles $\theta$, $\varpi$ relativement à cette molecule ; la distance de la molecule au point attiré sera $\sqrt{(r^2 - 2rR\mu + R^2)}$ en faisant pour abreger
\[ \mu = \cos\theta \cos\theta' + \sin\theta \sin\theta' \cos(\varpi - \varpi'), \]
la masse de la molecule sera $R^2 dR\, d\theta'\, d\varpi' \sin\theta'$, \struck{ou \uncertain{$R^2 dR$} \ill{}}, et si on designe par $V$ la somme des molecules divisées par leur distances au point attiré, on aura \struck{\uncertain{$V = \int$}}
\[ V = \int^{3} \frac{R^2 dR\, d\varpi'\, d\theta' \sin\theta'}{\sqrt{(r^2 - 2rR\mu + R^2)}} \]
\note{« $\sqrt{\ }$ » : le radical est écrit comme un grand \emph{V} suivi d'une parenthèse, sans barre ; il est rendu par $\sqrt{(\ )}$, la parenthèse gardée. « $\int^{3}$ » : un 3 est écrit en exposant après le signe d'intégration, pour l'intégrale triple. Une grosse tache d'encre couvre le signe $=$ de la formule détachée, qui se lit dessous. La ligne biffée après « $\sin\theta'$, » commence par « ou » et par « $R^2 dR$ » ; le reste, sous le trait, ne se lit pas. Après « on aura », un début de formule biffé, où l'on croit voir « $V = \int$ ».}

l'integrale relative à $R$ devant être prise depuis $R = 0$ jusqu'à la valeur de $R$ à la surface du sphéroïde, l'integrale relative à $\varpi'$, devant être prise depuis $\varpi' = 0$ jusqu'à $\varpi' = 360^\circ$, et celle qui est relative à $\theta'$ \uncertain{ou à $\mu$} depuis $\theta' = 0$ jusqu'à $\theta' = 180^\circ$ ; \struck{et que \ill{} depuis \ill{} jusqu'à \ill{}} La valeur de $V$ étant ainsi determinée on aura $\frac{dV}{dr}$, $\frac{dV}{r d\theta}$, \struck{\ill{}} $\frac{dV}{r d\varpi \sin\theta}$ pour les trois attractions du sphéroïde dans le sens du rayon $r$, \struck{dans} et perpendiculairement à ce rayon dans le plan $r$ et $R$, et dans un plan perpendiculaire à celui-ci.
\note{« ou à $\mu$ » est écrit dans la marge droite et appelé par un signe d'insertion après « $\theta'$ » ; le dernier signe a la forme du $\mu$ de la formule plus haut, mais la lecture reste douteuse. Le passage biffé après « $180^\circ$ ; » tient une ligne et demie ; seuls « et que », « depuis » et « jusqu'à » s'y lisent. Le $d$ de ces différentielles est le $d$ bouclé de la main, le même que dans « dans » ou « de » ; il est rendu par $d$. Entre la deuxième et la troisième expression, une fraction commencée puis biffée à traits serrés, illisible. « leur distances » : ainsi, sans \emph{s}. Le texte s'arrête sur « celui-ci. » vers les deux tiers de la page ; le reste est blanc, et le verso (vue 20) aussi.}

\end{document}
